\section{Torsion modules}
\label{section-torsion}

\noindent
In this section ``torsion modules'' will refer to modules supported
on a given closed subset $V(I)$ of an affine scheme $\Spec(R)$.
This is different, but analogous to, the notion of a torsion module
over a domain (Definition \ref{definition-torsion}).

\begin{definition}
\label{definition-f-power-torsion}
Let $R$ be a ring. Let $M$ be an $R$-module.
\begin{enumerate}
\item Let $I \subset R$ be an ideal. We say $M$ is an
{\it $I$-power torsion module} if for every $m \in M$ there exists an $n > 0$
such that $I^n m = 0$.
\item Let $f \in R$. We say $M$ is
{\it an $f$-power torsion module} if for each
$m \in M$, there exists an $n > 0$ such that $f^n m = 0$.
\end{enumerate}
\end{definition}

\noindent
Thus an $f$-power torsion module is the same thing as an $I$-power torsion
module for $I = (f)$. We will use the notation
$$
M[I^n] = \{m \in M \mid I^nm = 0\}
$$
and
$$
M[I^\infty] = \bigcup M[I^n]
$$
for an $R$-module $M$. Thus $M$ is $I$-power torsion if and only if
$M = M[I^\infty]$ if and only if $M = \bigcup M[I^n]$.

\begin{lemma}
\label{lemma-I-power-torsion-presentation}
Let $R$ be a ring.
Let $I$ be an ideal of $R$.
Let $M$ be an $I$-power torsion module.
Then $M$ admits a resolution
$$
\ldots \to K_2 \to K_1 \to K_0 \to M \to 0
$$
with each $K_i$ a direct sum of copies of $R/I^n$ for $n$ variable.
\end{lemma}

\begin{proof}
There is a canonical surjection
$$
\oplus_{m \in M} R/I^{n_m} \to M \to 0
$$
where $n_m$ is the smallest positive integer such that $I^{n_m} \cdot m = 0$.
The kernel of the preceding surjection is also an $I$-power torsion module.
Proceeding inductively, we construct the desired resolution of $M$.
\end{proof}

\begin{lemma}
\label{lemma-torsion-free}
Let $R$ be a ring. Let $I$ be an ideal of $R$.
For any $R$-module $M$ set $M[I^n] = \{m \in M \mid I^nm = 0\}$.
If $I$ is finitely generated then the following are equivalent
\begin{enumerate}
\item $M[I] = 0$,
\item $M[I^n] = 0$ for all $n \geq 1$, and
\item if $I = (f_1, \ldots, f_t)$, then the map
$M \to \bigoplus M_{f_i}$ is injective.
\end{enumerate}
\end{lemma}

\begin{proof}
This follows from
Algebra, Lemma \ref{algebra-lemma-when-injective-covering}.
\end{proof}

\begin{lemma}
\label{lemma-divide-by-torsion}
Let $R$ be a ring. Let $I$ be a finitely generated ideal of $R$.
\begin{enumerate}
\item For any $R$-module $M$ we have $(M/M[I^\infty])[I] = 0$.
\item An extension of $I$-power torsion modules is $I$-power torsion.
\end{enumerate}
\end{lemma}

\begin{proof}
Let $m \in M$. If $m$ maps to an element of $(M/M[I^\infty])[I]$
then $Im \subset M[I^\infty]$.
Write $I = (f_1, \ldots, f_t)$. Then we see that
$f_i m \in M[I^\infty]$, i.e., $I^{n_i}f_i m = 0$ for some $n_i > 0$.
Thus we see that $I^Nm = 0$ with $N = \sum n_i + 2$.
Hence $m$ maps to zero in $(M/M[I^\infty])$ which proves the
first statement of the lemma.

\medskip\noindent
For the second, suppose that $0 \to M' \to M \to M'' \to 0$ is a short
exact sequence of modules with $M'$ and $M''$ both $I$-power torsion
modules. Then $M[I^\infty] \supset M'$ and hence $M/M[I^\infty]$ is a
quotient of $M''$ and therefore $I$-power torsion. Combined with
the first statement and Lemma \ref{lemma-torsion-free} this implies
that it is zero.
\end{proof}

\begin{lemma}
\label{lemma-I-power-torsion}
Let $I$ be a finitely generated ideal of a ring $R$.
The $I$-power torsion modules form a Serre subcategory of
the abelian category $\text{Mod}_R$, see
Homology, Definition \ref{homology-definition-serre-subcategory}.
\end{lemma}

\begin{proof}
It is clear that a submodule and a quotient module of an $I$-power
torsion module is $I$-power torsion. Moreover, the extension
of two $I$-power torsion modules is $I$-power torsion by
Lemma \ref{lemma-divide-by-torsion}.
Hence the statement of the lemma by
Homology, Lemma \ref{homology-lemma-characterize-serre-subcategory}.
\end{proof}

\begin{lemma}
\label{lemma-local-cohomology-closed}
Let $R$ be a ring and let $I \subset R$ be a finitely generated ideal.
The subcategory $I^\infty\text{-torsion} \subset \text{Mod}_R$
depends only on the closed subset $Z = V(I) \subset \Spec(R)$.
In fact, an $R$-module $M$ is $I$-power torsion if and only if its
support is contained in $Z$.
\end{lemma}

\begin{proof}
Let $M$ be an $R$-module. Let $x \in M$. If $x \in M[I^\infty]$, then $x$
maps to zero in $M_f$ for all $f \in I$. Hence $x$ maps to zero in
$M_\mathfrak p$ for all $\mathfrak p \not \supset I$. Conversely, if $x$
maps to zero in $M_\mathfrak p$ for all $\mathfrak p \not \supset I$,
then $x$ maps to zero in $M_f$ for all $f \in I$. Hence if
$I = (f_1, \ldots, f_r)$, then $f_i^{n_i}x = 0$ for some $n_i \geq 1$.
It follows that $x \in M[I^{\sum n_i}]$. Thus $M[I^\infty]$ is
the kernel of $M \to \prod_{\mathfrak p \not \in Z} M_\mathfrak p$.
The second statement of the lemma follows and it implies the first.
\end{proof}

\noindent
The next two lemmas should probably go somewhere else.

\begin{lemma}
\label{lemma-bounded-above-mod-I}
Let $R$ be a ring. Let $I \subset R$ be an ideal. Let $K$ be an object of
$D(R)$ such that $H^i(K \otimes_R^\mathbf{L} R/I) = 0$ for $i > 0$. Then
\begin{enumerate}
\item $H^i(K \otimes_R^\mathbf{L} R/I^n) = 0$ for all $n \geq 1$
and $i > 0$,
\item $H^i(K \otimes_R^\mathbf{L} N) = 0$ for any $I$-power torsion
$R$-module $N$ and $i > 0$, and
\item for any $M \in D^b(R)$ whose cohomology modules $H^i(M)$
are $I$-power torsion and $0$ for $i > 0$ we have
$H^i(K \otimes_R^\mathbf{L} M) = 0$ for $i > 0$.
\end{enumerate}
\end{lemma}

\begin{proof}
Proof of (2). We can write $N = \bigcup N[I^n]$. We have
$K \otimes_R^\mathbf{L} N = \text{hocolim}_n K \otimes_R^\mathbf{L} N[I^n]$
as tensor products commute with colimits (details omitted; hint: represent
$K$ by a K-flat complex and compute directly).
Hence we may assume $N$ is annihilated by $I^n$.
Consider the $R$-algebra $R' = R/I^n \oplus N$
where $N$ is an ideal of square zero. It suffices to show that the object
$K' = K \otimes_R^\mathbf{L} R'$ of $D(R')$ has vanishing cohomology
in positive degrees.
We have a surjection $R' \to R/I$ of $R$-algebras whose kernel $J$
is nilpotent (any product of $n$ elements in the kernel is zero).
We have
$$
K \otimes_R^\mathbf{L} R/I =
(K \otimes_R^\mathbf{L} R') \otimes_{R'}^\mathbf{L} R/I =
K' \otimes_{R'}^\mathbf{L} R/I
$$
By assumption the complex $K \otimes_R^\mathbf{L} R/I$ has tor-amplitude
in $[-\infty, 0]$. Thus the conclusion by
Lemma \ref{lemma-check-tor-dimension-modulo-nilpotent}.

\medskip\noindent
Part (1) follows trivially from part (2).
Part (3) follows from part (2), induction on the number of nonzero
cohomology modules of $M$, and the distinguished triangles of truncation
from Derived Categories, Remark
\ref{derived-remark-truncation-distinguished-triangle}. Details omitted.
\end{proof}

\begin{lemma}
\label{lemma-derived-vanishing-mod-I}
Let $R$ be a ring. Let $I \subset R$ be an ideal. Let $K$ be an object of
$D(R)$ such that $K \otimes_R^\mathbf{L} R/I = 0$ in $D(R)$. Then
\begin{enumerate}
\item $K \otimes_R^\mathbf{L} R/I^n = 0$ for all $n \geq 1$,
\item $K \otimes_R^\mathbf{L} N = 0$ for any $I$-power torsion
$R$-module $N$,
\item $K \otimes_R^\mathbf{L} M = 0$ for any $M \in D^b(R)$ whose
cohomology modules are $I$-power torsion.
\end{enumerate}
\end{lemma}

\begin{proof}
Consequence of Lemma \ref{lemma-bounded-above-mod-I} (some details
omitted).
\end{proof}

\begin{lemma}
\label{lemma-torsion-completion}
\begin{reference}
Slight generalization of \cite[Lemme~1]{Beauville-Laszlo}.
\end{reference}
Let $R \to R'$ be a ring map. Let $I \subset R$ be an ideal such that
$R/I^n \to R'/I^nR'$ is an isomorphism for $n > 0$.
For any $I$-power torsion $R$-module $M$ the map $M \to M \otimes_R R'$
is an isomorphism.
For example, if $I$ is finitely generated and $R^\wedge$ is the completion
of $R$ with respect to $I$, then we have $M \cong M \otimes_R R^\wedge$.
\end{lemma}

\begin{proof}
If $M$ is annihilated by $I^n$, then
$$
M \otimes_R R' \cong
M \otimes_{R/I^n} R'/I^n R' \cong
M \otimes_{R/I^n} R/I^n \cong M.
$$
If $M$ is $I$-power torsion, then $M = \bigcup M[I^n]$.
Since tensor products commute with direct limits
(Algebra, Lemma \ref{algebra-lemma-tensor-products-commute-with-limits}),
we obtain the desired isomorphism. The last statement is a special
case of the first statement by Algebra, Lemma
\ref{algebra-lemma-hathat-finitely-generated}.
\end{proof}







